\chapter{Rays, filtrations, and Legendre transforms}
\index{geodesic ray!filtration correspondence}\index{filtration!geodesic ray}\index{Legendre transform!comparison}
\label{chap:rays-filtrations}

A test configuration has four useful shadows.
It defines a filtration of the section ring, a Fubini--Study metric on the
trivially valued analytification, a curve of complex singularity types, and
a weak geodesic ray.
For algebraic input these constructions commute.
For general filtrations only some arrows remain, because finite generation
is exactly what reconstructs a test configuration.

This chapter fixes one sign convention and then proves the commutative
diagram.
It also audits the asymptotic-slope literature against the published 2022
erratum to \cite{BHJ19}.
No result depending on the flawed Theorem~5.6 of that paper is used.

Throughout the algebraic part, \(X\) is a projective complex manifold of
dimension \(n\), \(L\) is ample, and
\[
                       R(X,L)=\bigoplus_{m\geq0}R_m,
 \qquad R_m=H^0(X,mL).
\tag{65.1}\label{eq:rf-section-ring}
\]
We pass to sufficiently divisible \(m\) whenever \(L\) or a test
configuration polarization is rational.

\section{From a test configuration to a filtration}

Let $X$ be a connected projective complex manifold, let $L$ be ample, and let $(\mathcal X,\mathcal L)$ be a relatively ample test configuration for $(X,L)$, with parameter $\varpi$ on $\mathbb A^1$. The total space is not assumed normal in this construction. Work in degrees $m$ for which $mL$ and $m\mathcal L$ are line bundles. The punctured-product trivialization associates to $s\in R_m$ an invariant section $\bar s$ over $\mathcal X\setminus\mathcal X_0$.

\begin{definition}\label{def:rf-tc-filtration}
For $\lambda\in\mathbb R$, set
\[
 \mathcal F^\lambda R_m
 :=\left\{s\in R_m\ \middle|\
       \varpi^{-\lceil\lambda\rceil}\bar s
       \in H^0(\mathcal X,m\mathcal L)\right\}.
\tag{65.2}\label{eq:rf-tc-filtration}
\]
This is the filtration of the actual model, before normalization. For normal models it agrees with \Cref{def:ch33-tc-filtration}, whose convention first takes integral closure in the generic fibre.
\end{definition}

The ceiling makes the filtration left-continuous. For $0\ne s\in R_m$, its additive weight is
\[
               \chi_m(s)=\sup\{\lambda\mid s\in\mathcal F^\lambda R_m\}.
\tag{65.3}\label{eq:rf-additive-weight}
\]
This is an integer and the supremum is attained. Our convention assigns weight $\lambda$ to $s\varpi^{-\lambda}$. A character twist of weight $c$ adds $mc$ to $\chi_m$; rational $c$ is allowed after clearing denominators in the polarization. A real shift also defines a filtration, but need not be an algebraic character twist.

\begin{proposition}\label{prop:rf-filtration-properties}
The filtration \eqref{eq:rf-tc-filtration} is decreasing, exhaustive, separated, left-continuous, multiplicative, and linearly bounded. Its Rees algebra is finitely generated over $\mathbb C[\varpi]$ after clearing the polarization denominators.
\end{proposition}

\begin{proof}
For each admissible $m$, the module $H^0(\mathcal X,m\mathcal L)$ is a finite torsion-free $\mathbb C[\varpi]$-module whose localization is $R_m[\varpi,\varpi^{-1}]$. It is therefore a full lattice. An equivariant basis of this lattice gives finitely many integer jumps, proving exhaustivity, separation and attainment of \eqref{eq:rf-additive-weight}. Multiplication by $\varpi$ proves decreasingness. The ceiling convention gives left-continuity.

The product of the extending sections attached to $s\in\mathcal F^\lambda R_m$ and $s'\in\mathcal F^\mu R_{m'}$ has pole exponent $\lceil\lambda\rceil+\lceil\mu\rceil$. Since this is at least $\lceil\lambda+\mu\rceil$, multiplication by the corresponding nonnegative power of $\varpi$ gives
\[
 \mathcal F^\lambda R_m\cdot\mathcal F^\mu R_{m'}
       \subset\mathcal F^{\lambda+\mu}R_{m+m'}.
\tag{65.4}\label{eq:rf-multiplicativity}
\]
The weight decomposition identifies the Rees algebra with the relative section algebra $\bigoplus_m H^0(\mathcal X,m\mathcal L)$. Relative ampleness makes it finitely generated. Choose homogeneous equivariant generators $s_j\varpi^{-\lambda_j}$, of positive section degrees $d_j$, together with $\varpi$ in degree zero. Their sections $s_j$ generate the generic section ring. Thus every section of degree $m$ is a sum of monomials with weight at least $m\min_j(\lambda_j/d_j)$. Conversely, a homogeneous expression in the Rees algebra has weight at most $m\max_j(\lambda_j/d_j)$, since a power of $\varpi$ lowers weight. Any constants $C_-<\min_j(\lambda_j/d_j)$ and $C_+>\max_j(\lambda_j/d_j)$ satisfy
\[
 \mathcal F^{mC_-}R_m=R_m,\qquad \mathcal F^{mC_+}R_m=0
 \quad(m>0).
\tag{65.5}\label{eq:rf-linear-bounds}
\]
These constants are independent of $m$. The lattice and Rees identifications are those of \cite[Section~1.2 and Proposition~2.15]{BHJ17}.
\end{proof}

For the converse, let $\mathcal F$ be a multiplicative, exhaustive, separated, left-continuous filtration with the bounds \eqref{eq:rf-linear-bounds}. In the Rees statements, ``rational'' has the convention of \Cref{def:ch33-fg-filtration}: after a Veronese restriction the jumping numbers are integral. Put $\mathcal F^\lambda R_0=\mathbb C$ for $\lambda\leq0$, and zero otherwise. Its Rees algebra is
\[
 A=\mathcal R(\mathcal F)
 :=\bigoplus_{m\geq0}\bigoplus_{\lambda\in\mathbb Z}
       \mathcal F^\lambda R_m\,\varpi^{-\lambda}
       \subset R(X,L)[\varpi,\varpi^{-1}].
\tag{65.6}\label{eq:rf-rees-algebra}
\]
The $\Proj$ grading is the section degree $m$; the second grading records the action. In particular $A_0=\mathbb C[\varpi]$.

\begin{theorem}\label{thm:rf-rees-reconstruction}
Assume $A$ is finitely generated. Choose $r$ so that $A^{(r)}=\bigoplus_{k\geq0}A_{rk}$ is generated in degree one. Then
\[
 \mathcal Y=\Proj_{\mathbb C[\varpi]}A^{(r)},\qquad
 \mathcal M=r^{-1}\mathcal O_{\mathcal Y}(1)
\]
is a relatively ample test configuration for $(X,L)$. Its normalization $\nu:\mathcal Y^\nu\to\mathcal Y$, with polarization $\nu^*\mathcal M$, is a normal relatively ample test configuration. The filtration obtained from $\mathcal Y$ is the section saturation of $A$; normalization can enlarge it further. Starting from a normal relatively ample configuration recovers that configuration.

With polarization and linearization retained, the resulting normal correspondence is
\[
 \mathsf{TC}^{\mathrm{n,ra}}
       \xleftrightarrow{\ \sim\ }
 \mathsf{Fil}^{\mathrm{fg},\mathbb Q}/\!\sim_{\mathrm{sat,ic}}.
\tag{65.7}\label{eq:rf-rees-equivalence}
\]
Here the left side consists of normal relatively ample test configurations up to equivariant polarized isomorphism. On the right, the equivalence relation is equality after passing to saturated integral closures of Rees algebras and a common Veronese restriction. Section saturation means replacing $A_{rk}$ by $H^0(\mathcal Y,\mathcal O_{\mathcal Y}(k))$ inside the common generic section space.
\end{theorem}

\begin{proof}
Generation in degree one gives projectivity and the relatively ample $\mathcal O(1)$. Every affine $\Proj$ chart is a subring of the common function field, so it is torsion-free over the principal ideal domain $\mathbb C[\varpi]$ and hence flat. No torsion needs to be discarded. Exhaustivity gives
\[
 A^{(r)}[\varpi^{-1}]
       =R(X,rL)[\varpi,\varpi^{-1}],
\]
which identifies the punctured family with $X\times\mathbb G_m$ and the polarization with $L$. The weight grading gives the equivariant structure.

The natural maps $A_{rk}\to H^0(\mathcal Y,\mathcal O(k))$ are isomorphisms for sufficiently large $k$. This is relative section-ring reconstruction; after a further Veronese restriction it gives the ample Rees correspondence of \cite[Proposition~2.15]{BHJ17}. The normalization is finite, since $\mathcal Y$ is of finite type over $\mathbb C$, and it is an isomorphism over the normal punctured product. Finite pullback preserves ampleness. Its affine charts remain torsion-free over $\mathbb C[\varpi]$, so the normalized family is also flat. Integral closure commutes with localization; taking the degree-zero parts on the homogeneous charts therefore gives the normalization of $\Proj A^{(r)}$. Taking all sections of the pulled-back polarization then gives the section saturation of this integral closure in the generic section ring.

For a normal relatively ample configuration the full section algebra reconstructs the original polarized scheme, and no normalization is needed. More generally, pulling back to a common normal domination preserves these sections by the projection formula and $\mu_*\mathcal O=\mathcal O$; see \cite[Lemma~2.13]{BHJ17}. Conversely equality of the saturated section algebras recovers the same polarized $\Proj$. The uniqueness of the normal ample representative in a semiample equivalence class is \cite[Proposition~2.30]{BJ22}. This proves \eqref{eq:rf-rees-equivalence}. The quotient on the filtration side is essential.
\end{proof}

\begin{samepage}
\begin{warning}{Infinite generation}\label{warn:rf-infinite-generation}
A linearly bounded multiplicative filtration need not be finitely generated. It still defines asymptotic metrics and a DH measure, but \Cref{thm:rf-rees-reconstruction} does not produce a finite-type test configuration from it.
\end{warning}
\end{samepage}

\section{The Fubini--Study metric}

Let $X$ be a connected projective complex manifold with ample $L$, and let $\mathcal F$ be a linearly bounded multiplicative filtration of its full section ring. For a semivaluation $v\in X^{\An}$, evaluate a section $s\in R_m$ in a local frame of $mL$ at the center of $v$. A change of frame multiplies by a unit of value zero, so $v(s)\in[0,+\infty]$ is well defined. A section vanishing on the support of $v$ has value $+\infty$.

\begin{definition}\label{def:rf-fs-metric}
The Fubini--Study function of $\mathcal F$ is
\[
 \phi_{\mathcal F}(v)
 :=\sup_{m\geq1}\ \sup_{0\ne s\in R_m}
             \frac{\chi_m(s)-v(s)}{m}.
\tag{65.8}\label{eq:rf-fs-metric}
\]
A term with $v(s)=+\infty$ is $-\infty$.
\end{definition}

There is no subtraction of two infinite values, since $\chi_m(s)$ is finite. In fact the whole supremum is finite even at proper-support semivaluations. The upper linear bound and $v(s)\geq0$ give $\phi_{\mathcal F}\leq C_+$. For a sufficiently divisible globally generated degree $m$, a section not vanishing at the center has value zero and weight at least $mC_-$. Thus $C_-\leq\phi_{\mathcal F}$ everywhere. Individual filtered pieces may have base loci; the full series has no common base locus in large divisible degrees.

\begin{proposition}\label{prop:rf-fg-pl}
If $\mathcal F$ is finitely generated and rational in the Rees convention above, there is a degree $r$ and a basis $s_1,\ldots,s_N$ of $R_r$ adapted to its filtration such that
\[
 \phi_{\mathcal F}(v)
       =\frac1r\max_{1\leq j\leq N}
                     \{\chi_r(s_j)-v(s_j)\}.
\tag{65.9}\label{eq:rf-finite-fs}
\]
It is the rational semipositive PL function of the Rees test configuration and of its normalization.
\end{proposition}

\begin{proof}
Restriction to a Veronese ring does not change \eqref{eq:rf-fs-metric}: a section $s$ in an omitted degree may be replaced by a power $s^q$ in an included degree, with $\chi(s^q)\geq q\chi(s)$ and $v(s^q)=qv(s)$. Choose $r$ so that $A^{(r)}$ is generated by its degree-one lattice, and take an equivariant lattice basis $s_j\varpi^{-\chi_r(s_j)}$. Each filtered section is a sum of monomials in these generators multiplied by nonnegative powers of $\varpi$. The quotient of each monomial in \eqref{eq:rf-fs-metric} is bounded by the maximum in \eqref{eq:rf-finite-fs}; the ultrametric inequality gives the same bound for its sum. The reverse bound holds because each generator occurs in the supremum. The basis generates $rL$ after increasing $r$, so the maximum is finite on all of $X^{\An}$.

On a common model dominating the Rees configuration and the trivial family, the generating sections give this same maximum for the vertical divisor of the polarization. Pullback to the normalization does not change that model function. This is the finite-section realization of \cite[Theorem~2.31 and Corollary~2.32]{BJ22}. The direct algebraic invariance, including proper-support points, is proved next.
\end{proof}

\begin{lemma}\label{lem:rf-integral-closure}
For a coherent ideal $\mathfrak a$ on $X$ and a semivaluation $v$, one has $v(\mathfrak a)=v(\overline{\mathfrak a})$. If $\overline{\mathcal F}$ is defined by the integral closure $\overline A$ of a finitely generated Rees algebra $A$ in $R(X,L)[\varpi,\varpi^{-1}]$, then $\phi_{\overline{\mathcal F}}=\phi_{\mathcal F}$. Section saturation also preserves this function. These assertions do not say that the individual filtered pieces are unchanged.
\end{lemma}

\begin{proof}
Inclusion gives $v(\overline{\mathfrak a})\leq v(\mathfrak a)$. If $f\in\overline{\mathfrak a}$ and $v(f)<v(\mathfrak a)$, an integral equation $f^k+a_1f^{k-1}+\cdots+a_k=0$, with $a_j\in\mathfrak a^j$, has a unique term of least value, namely $f^k$. This contradicts the ultrametric inequality. The argument also applies when $v(\mathfrak a)=+\infty$, since the supposed value $v(f)$ is then finite. This proves the ideal assertion.

For the graded assertion, fix $v$ and put $b=\phi_{\mathcal F}(v)\in\mathbb R$. Choose a local frame of $L$ at its center; after clearing denominators this may be done with an honest line bundle. Introduce a variable $z$ recording section degree. The Gauss extension $W$ of $v$ with $W(z)=b$ and $W(\varpi)=1$ satisfies
\[
 W(s z^m\varpi^{-\lambda})=v(s)+mb-\lambda\geq0
 \quad\text{for }s\in\mathcal F^\lambda R_m.
\]
Here a section is represented by its coefficient in the chosen frame. Gauss extension is valid for a semivaluation by first quotienting by its support and then extending the resulting valuation. Thus $W$ is nonnegative on $A$. If a homogeneous $h$ integral over $A$ had $W(h)<0$, its monic integral equation would again have a unique least-value term. Hence $W(h)\geq0$. Taking $h=s z^m\varpi^{-\lambda}$ shows $\lambda-v(s)\leq mb$ for every new filtered section with finite value; sections of infinite value contribute $-\infty$. Therefore $\phi_{\overline{\mathcal F}}\leq\phi_{\mathcal F}$, and inclusion gives equality.

Finally, section saturation agrees with the original graded algebra in all sufficiently large degrees. A section in any added positive degree has a power in that range, and the same power argument used in \Cref{prop:rf-fg-pl} bounds its Fubini--Study contribution. Thus section saturation also preserves the function. This proves invariance of the supremum without asserting invariance of any fixed-degree base ideal.
\end{proof}

\begin{example}\label{ex:rf-normalization-changes-weights}
Take $X=\mathbb P^1$ and $L=\mathcal O(1)$, with section ring $\mathbb C[x,y]$. The Rees algebra
\[
 A=\mathbb C[\varpi,x,y,\varpi^{-2}x^2]
      \subset\mathbb C[x,y,\varpi,\varpi^{-1}]
\]
defines a multiplicative filtration with monomial weights
\[
 \chi_m(x^a y^{m-a})=2\lfloor a/2\rfloor,
 \qquad 0\leq a\leq m.
\tag{65.9a}\label{eq:rf-nonnormal-weights}
\]
The filtered pieces are spanned by the monomials satisfying the weight threshold. They are linearly bounded since $0\leq2\lfloor a/2\rfloor\leq m$. Put $u=\varpi^{-1}x$. Then $u$ is integral over $A$, since $u^2=\varpi^{-2}x^2$, and
\[
                    \overline A=\mathbb C[\varpi,u,y].
\]
Indeed this polynomial ring is integral over $A$, is normal, and has the same fraction field. The new monomial weight is $\overline\chi_m(x^a y^{m-a})=a$. In particular $\mathcal F^1R_1=0$, whereas $\overline{\mathcal F}^{1}R_1=\mathbb Cx$. The corresponding base ideals are $0$ and the ideal of $(x=0)$, whose values differ wherever $v(x)$ is finite. Nevertheless
\[
 \phi_{\mathcal F}(v)=\phi_{\overline{\mathcal F}}(v)
                         =\max\{-v(y),\,1-v(x)\}.
\tag{65.9b}\label{eq:rf-normalization-same-fs}
\]
For the original filtration the two terms are already realized by $y$ and $x^2$ in degrees one and two. The remaining monomials give convex combinations bounded by this maximum. For the normalized filtration the terms are realized by $y$ and $x$. This also verifies the equality at proper-support semivaluations.

The same example tests normalized base change. Start with
\[
 A_0=\mathbb C[t,x,y,t^{-1}x^2],\qquad
       \chi^0_m(x^a y^{m-a})=\lfloor a/2\rfloor.
\]
This ring is normal: it is the intersection of $\mathbb C[t,t^{-1},x,y]$ with the valuation ring of the weight valuation $w(t)=2$, $w(x)=1$, $w(y)=0$. Both rings are integrally closed. After $t=\varpi^2$, the unnormalized base change has algebra $A$ and weights $2\chi^0_m$. Its normalization has weights $a$, which are strictly larger for odd $a$. Such monomials occur in arbitrarily large divisible section degrees.

These weights occur in the full section lattices of actual test configurations. Indeed, the second Veronese algebras have presentations
\[
 \begin{aligned}
 A_0^{(2)}&=\mathbb C[t,U,V,W]/(V^2-tUW),\\
 A^{(2)}&=\mathbb C[\varpi,U,V,W]/(V^2-\varpi^2UW),
 \end{aligned}
\]
where $U=y^2$, $V=xy$, and $W=t^{-1}x^2$ or $\varpi^{-2}x^2$, respectively. Give $U,V,W$ degree one. The relative $\Proj$ schemes are plane conics with polarization $\mathcal O(1)/2$. Their hypersurface exact sequences identify the displayed graded pieces with the full spaces of sections in every nonnegative degree: the obstruction group is $H^1(\mathbb P^2_B,\mathcal O(k-2))=0$ for $B=\mathbb C[t]$ or $\mathbb C[\varpi]$. Thus the discrepancy is not removed by section saturation. Normalized base change preserves the expected scaling of the metric, but its finite-level weights need not equal the base-change degree times the old weights.
\end{example}

\section{The finite-level algebra behind the metric}
\label{sec:rf-finite-level-algebra}

The passage from a filtration to a metric is easiest to audit before taking
the supremum over \(m\).  For each \(m\) and \(\lambda\), let
\(\mathfrak a_{m,\lambda}\) be the base ideal of
\(\mathcal F^\lambda R_m\), namely the image of
\[
 \mathcal F^\lambda R_m\otimes\mathcal O_X(-mL)
                  \longrightarrow\mathcal O_X.
\tag{65.10}\label{eq:rf-filtered-base-ideal}
\]
Multiplicativity of \(\mathcal F\) gives
\[
 \mathfrak a_{m,\lambda}\mathfrak a_{m',\mu}
       \subset
 \mathfrak a_{m+m',\lambda+\mu}.
\tag{65.11}\label{eq:rf-base-ideal-multiplicativity}
\]
If \(s_1,\ldots,s_N\) generate the filtered piece and \(v\) is a
semivaluation, then
\[
 v(\mathfrak a_{m,\lambda})
       =\min_jv(s_j),
 \qquad
 \phi_{m,\lambda}(v)
       =\frac{\lambda-v(\mathfrak a_{m,\lambda})}{m}.
\tag{65.12}\label{eq:rf-base-ideal-fs-level}
\]
Taking the supremum over all nonzero levels recovers
\eqref{eq:rf-fs-metric}.  Thus the filtration-to-FS arrow factors as
\[
 \mathcal F
 \longmapsto
 \{\mathfrak a_{m,\lambda}\}_{m,\lambda}
 \longmapsto
 \left(v\mapsto
   \sup_{m,\lambda}
   \frac{\lambda-v(\mathfrak a_{m,\lambda})}{m}\right).
\tag{65.13}\label{eq:rf-filtration-ideal-metric-factorization}
\]
This factorization is the common algebraic input for the BJ metric and the
complex partial Bergman obstacles: the non-Archimedean side evaluates the
base ideals by \(v\), while the complex side evaluates their sections by
logarithmic squared norms and then takes maximal envelopes.

For each fixed ideal, integral closure preserves its value by \Cref{lem:rf-integral-closure}. Normalizing the entire Rees algebra is a different operation: it may enlarge a filtered piece by using an integral relation involving higher degrees. Hence the individual values $v(\mathfrak a_{m,\lambda})$ may change. What remains invariant is the supremum in \eqref{eq:rf-filtration-ideal-metric-factorization}, by the graded argument of \Cref{lem:rf-integral-closure}. The difference is explicit in \Cref{ex:rf-normalization-changes-weights}.

\subsection{A three-check reconstruction audit}

Given a proposed finitely generated filtration/test-configuration
correspondence, one can verify it in three finite computations.

\begin{enumerate}
\item Localizing the Rees algebra at the parameter must give
\(R(X,L)[\varpi,\varpi^{-1}]\).  Otherwise the generic fibre is not the
punctured product.

\item The degree-\(m\) weight of a section must equal its maximal allowed
pole order along the central fibre, with the sign in
\eqref{eq:rf-tc-filtration}.  Otherwise the FS maximum is reflected.

\item Twisting the linearization by \(c\) must add \(mc\) in degree \(m\),
translate the FS metric by \(c\), and translate the DH measure by \(c\).
If only two of these changes occur, the linearization was lost in one of the
arrows.
\end{enumerate}

\paragraph{Source boundary for the finite-level factorization.}
The Rees correspondence is \cite[Proposition~2.15]{BHJ17}; normal ample representatives and their Fubini--Study functions are \cite[Proposition~2.30 and Theorem~2.31]{BJ22}.  Equations
\eqref{eq:rf-filtered-base-ideal}--
\eqref{eq:rf-filtration-ideal-metric-factorization} merely expand the
finite-generator calculation used in \Cref{prop:rf-fg-pl}; they do not
assert finite generation for a general filtration.

\section{Filtered sections and test curves}
\label{sec:rf-filtered-test-curves}

Let $X$ be a connected projective complex manifold of dimension $n$, let $L$ be ample, and let $h$ be a smooth positive Hermitian metric on $L$. Put $\theta=c_1(L,h)$ and use zero as the reference potential. Let $\mathcal F$ be a decreasing, left-continuous, multiplicative, linearly bounded filtration of the full section ring $R(X,L)$. For a rational polarization all degrees below are sufficiently divisible. Fix a smooth probability volume form $\mathrm dV$ and write
\[
 \|s\|_{2,m}^2=\int_X|s|_{h^m}^2\,\mathrm dV,
 \qquad \|s\|_{\infty,m}^2=\sup_X|s|_{h^m}^2.
\]
For $\tau\in\mathbb R$, the sup-norm partial Bergman obstacle is
\[
 B_{m,\tau}
 :=\uscsup\left\{\frac1m\log|s|_{h^m}^2\ \middle|\
 s\in\mathcal F^{m\tau}R_m,\ \|s\|_{\infty,m}\leq1\right\}.
\tag{65.14}\label{eq:rf-partial-bergman}
\]
It is $-\infty$ when the filtered space is zero. Let $B^{(2)}_{m,\tau}$ be the same extremum with $\|s\|_{2,m}\leq1$. For an $L^2$-orthonormal basis $s_{m,j}$ adapted to the filtration, with weights $\lambda_{m,j}$, the norm of the evaluation functional gives
\[
 B^{(2)}_{m,\tau}
   =\frac1m\log\sum_{\lambda_{m,j}\geq m\tau}|s_{m,j}|_{h^m}^2.
\]
The two normalizations agree asymptotically, with a bound uniform in the filtered subspace.

\begin{lemma}\label{lem:rf-bergman-norm-comparison}
There is a constant $C\geq1$, depending only on $X,L,h,\mathrm dV$, such that
\[
 \begin{gathered}
 \|s\|_{2,m}^2\leq\|s\|_{\infty,m}^2
          \leq C m^n\|s\|_{2,m}^2,\\
 B_{m,\tau}\leq B^{(2)}_{m,\tau}
       \leq B_{m,\tau}+\varepsilon_m,
 \qquad \varepsilon_m=\frac{\log C+n\log m}{m}.
 \end{gathered}
\tag{65.14a}\label{eq:rf-bergman-norm-comparison}
\]
The inequalities include zero spaces and common zero loci, using $-\infty+\varepsilon_m=-\infty$.
\end{lemma}

\begin{proof}
The first inequality uses the unit mass of $\mathrm dV$. For the reverse one choose local coordinates and a holomorphic frame at a point so that the first derivative of the metric weight vanishes there. On a ball of radius $m^{-1/2}$ the variation of $m$ times that weight is uniformly bounded. The holomorphic mean-value inequality on this ball bounds the squared value at its center by $C m^n$ times the global squared $L^2$-norm. Compactness makes $C$ uniform in the center; finitely many small degrees can be absorbed into $C$. The evaluation-functional formula and rescaling the two unit balls give the second line. This is the uniform norm comparison in the proof of \cite[Lemmas~7.4 and~7.6]{RWN14}.
\end{proof}

Put
\[
 g_\tau^{\mathcal F}:=\uscsup_{m\geq1}B_{m,\tau},
 \qquad
 c_{\mathcal F}:=\sup_{\substack{m\geq1\\R_m\ne0}}
           \max_{0\ne s\in R_m}\frac{\chi_m(s)}m.
\]
Linear boundedness makes $c_{\mathcal F}$ finite. Define the interior levels by
\[
 \Gamma_\tau^{\mathcal F}:=P_\theta[g_\tau^{\mathcal F}],
                 \qquad \tau<c_{\mathcal F},
\tag{65.15}\label{eq:rf-filtration-test-curve}
\]
and use the closed endpoint convention
\[
 \Gamma_{c_{\mathcal F}}^{\mathcal F}
       :=\inf_{\tau<c_{\mathcal F}}\Gamma_\tau^{\mathcal F},
 \qquad
 \Gamma_\tau^{\mathcal F}:=-\infty\quad(\tau>c_{\mathcal F}).
\tag{65.15a}\label{eq:rf-filtration-closed-endpoint}
\]
The endpoint is taken from the interior levels, not from a single extremal filtered piece.

\begin{theorem}\label{thm:rf-filtration-test-curve}
The family $\Gamma^{\mathcal F}$ is a bounded maximal test curve with terminal parameter $c_{\mathcal F}$ and reference potential zero. For $\tau<c_{\mathcal F}$ one has $g_\tau^{\mathcal F}=\Gamma_\tau^{\mathcal F}$. Its Legendre transform
\[
 \ell_t^{\mathcal F}:=\sup_{\tau\in\mathbb R}
             \bigl(\Gamma_\tau^{\mathcal F}+t\tau\bigr),\qquad t>0,
\tag{65.16}\label{eq:rf-filtration-ray}
\]
is a weak geodesic ray from zero. The supremum is already $\theta$-psh, and
\[
              C_-t\leq\ell_t^{\mathcal F}\leq c_{\mathcal F}t.
\]
Here $C_-$ is as in \eqref{eq:rf-linear-bounds}. The radial parameter for squared-norm potentials is $t=-\log|\zeta|^2$.
\end{theorem}

\begin{proof}
Powers of normalized sections give $B_{qm,\tau}\geq B_{m,\tau}$. Hence for every $k$ the regularized supremum over $m\geq k$ equals $g_\tau^{\mathcal F}$. Along factorial degrees the obstacles increase and recover the same regularized supremum. Combining this with \eqref{eq:rf-bergman-norm-comparison} gives
\[
 g_\tau^{\mathcal F}
 =\lim_{k\to\infty}\left(\sup_{m\geq k}B^{(2)}_{m,\tau}\right)^*.
\tag{65.16a}\label{eq:rf-partial-bergman-tail}
\]
This argument allows zero pieces in infinitely many degrees; it makes no assertion of pointwise convergence along every degree at such a level.

For concavity, take normalized sections in degrees $m,m'$ at parameters $\tau,\tau'$. If $a=p/q\in[0,1]\cap\mathbb Q$, their powers with exponents $pm'$ and $(q-p)m$ give a section in degree $qmm'$ at parameter $a\tau+(1-a)\tau'$. Its normalized logarithm is the corresponding convex combination. Taking suprema yields the concavity inequality outside the pluripolar exceptional sets of upper regularization; uniqueness of psh representatives extends it to all of $X$. Monotonicity in the parameter extends the inequality from rational to real $a$. This is the argument of \cite[Proposition~7.7]{RWN14}.

The definition of $c_{\mathcal F}$ gives a nonzero filtered section at each $\tau<c_{\mathcal F}$, so $g_\tau^{\mathcal F}\not\equiv-\infty$ and it has small unbounded locus. The analytic input is the interior maximality theorem \cite[Proposition~7.11 and Corollary~7.12]{RWN14}. Together with \cite[proof of Proposition~6.10, equation~(6.6)]{RWN14}, it says that the regularized supremum just constructed is fixed by the maximal envelope for every $\tau<c_{\mathcal F}$. Thus $g_\tau^{\mathcal F}=P_\theta[g_\tau^{\mathcal F}]$. This step uses the source's division and envelope argument; it does not follow just from multiplicativity.

For $\tau\leq C_-$ a globally generated degree gives a bounded lower obstacle, while $g_\tau^{\mathcal F}\leq0$. Its maximal envelope is therefore zero. Above $c_{\mathcal F}$ all filtered pieces vanish. Interior maximality gives supremum zero at each interior level. Their nested compact zero sets prevent the decreasing endpoint from collapsing identically to $-\infty$. By \Cref{lem:rays-closed-endpoint}, \eqref{eq:rf-filtration-closed-endpoint} gives the closed concave extension. The bounded RWN correspondence, \cite[Corollary~7.12]{RWN14} and \Cref{thm:rays-rwn}, now gives the weak geodesic. The absence of a further regularization is \Cref{prop:rays-legendre-well-defined}. Restricting the supremum to $[C_-,c_{\mathcal F}]$ proves the bounds and uniform convergence to zero as $t\downarrow0$.
\end{proof}

The raw value $g_{c_{\mathcal F}}^{\mathcal F}$ is at most the closed endpoint, and that endpoint is the decreasing limit of the interior levels. Consequently $g^{\mathcal F}$ and $\Gamma^{\mathcal F}$ have the same forward Legendre supremum. The inverse transform, however, recovers the closed curve:
\[
 \Gamma_\tau^{\mathcal F}
       =\inf_{t>0}(\ell_t^{\mathcal F}-t\tau).
\tag{65.17}\label{eq:rf-inverse-legendre}
\]
The value $-\infty$ is retained when the infimum diverges downward.

\begin{proposition}\label{prop:rf-legendre-biconjugacy}
On bounded maximal test curves with the closed endpoint convention and their weak geodesic rays, \eqref{eq:rf-filtration-ray} and \eqref{eq:rf-inverse-legendre} are inverse. Translation of filtration weights by $mc$ gives
\[
 \Gamma_\tau\longmapsto\Gamma_{\tau-c},
 \qquad \ell_t\longmapsto\ell_t+ct.
\tag{65.18}\label{eq:rf-translation}
\]
\end{proposition}

\begin{proof}
For every $x$, the profile $\tau\mapsto\Gamma_\tau(x)$ is closed, concave and decreasing, and is zero on a lower half-line. Its conjugate needs only slopes $t\geq0$; the value at $t=0$ is recovered by taking $t\downarrow0$. Fenchel--Moreau duality with the sign $+t\tau$ therefore gives the infimum over $t>0$. By \Cref{prop:rays-legendre-well-defined}, the forward supremum is already psh, so this scalar duality applies to the actual ray. For $\tau>c_{\mathcal F}$ the upper bound $\ell_t\leq c_{\mathcal F}t$ gives the required value $-\infty$. The endpoint is included by closedness. Substituting $\tau-c$ in the supremum proves \eqref{eq:rf-translation}, including the translated endpoint.
\end{proof}

\begin{example}\label{ex:rf-unattained-filtration-endpoint}
On $(\mathbb P^1,\mathcal O(1))$ with the Fubini--Study reference, give every nonzero section of degree $m\geq1$ weight $m-1$. This is multiplicative because $m+m'-1\geq(m-1)+(m'-1)$; its filtration is left-continuous and linearly bounded. Here $c_{\mathcal F}=1$, $g_\tau^{\mathcal F}=0$ for $\tau<1$, but $g_1^{\mathcal F}=-\infty$: no degree attains normalized weight one. The closed endpoint is $\Gamma_1^{\mathcal F}=0$ and the ray is $\ell_t=t$. Its inverse at one equals zero, so it cannot recover the raw value $g_1^{\mathcal F}$.

If instead every degree-$m$ section has weight $2\lfloor m/2\rfloor$, multiplicativity again holds. At $\tau=1$ all even-degree spaces are full and all odd-degree spaces vanish. The partial obstacles have no pointwise limit along all degrees: the even-degree sup-norm obstacles are zero for this reference, and the odd ones are $-\infty$. Indeed, at each point a unit linear form attains Fubini--Study norm one, and its $m$th power gives the full-series obstacle zero in degree $m$. Powers and tail suprema still give $g_1^{\mathcal F}=0$ and the same closed curve and ray. Thus left-continuity in the filtration index does not justify an unrestricted pointwise degree limit or supply the closed terminal value by itself.
\end{example}

\section{The factor-two firewall}
\label{sec:rf-factor-two-firewall}

The complex literature uses both norm weights and squared-norm potentials.
They must not share the same radial parameter without conversion.  Let
\[
       s=-\log|\zeta|,
       \qquad t=-\log|\zeta|^2=2s,
       \qquad \varphi_t=2\psi_s,
\tag{65.19}\label{eq:rf-factor-two-basic}
\]
where \(\psi_s=-\log|\cdot|\) is an ordinary norm weight and
\(\varphi_t=-\log|\cdot|^2\) is the book's squared-norm potential.
For a section of filtration weight \(\lambda\), the group action contributes
\[
        |\zeta|^{-2\lambda}
             =\rme^{2s\lambda}
             =\rme^{t\lambda}.
\tag{65.20}\label{eq:rf-factor-two-section-weight}
\]
This is exactly the exponential in
\eqref{eq:rf-bergman-ray}.  Replacing it by
\(\rme^{2t\lambda}\) would count the passage to squared norms twice.

The translation test makes the same point at the level of rays.  A shift
\(\chi_m\mapsto\chi_m+mc\) gives
\[
       \psi_s\longmapsto\psi_s+cs,
       \qquad
       \varphi_t\longmapsto\varphi_t+ct,
       \qquad t=2s.
\tag{65.21}\label{eq:rf-factor-two-translation}
\]
Hence the following three quotients encode the same slope:
\[
 \lim_{s\to\infty}\frac{\psi_s}{s},
 \qquad
 \lim_{s\to\infty}\frac{\varphi_{2s}}{2s},
 \qquad
 \lim_{t\to\infty}\frac{\varphi_t}{t}.
\tag{65.22}\label{eq:rf-factor-two-slope-quotients}
\]
A denominator \(s\) under a squared-norm potential doubles every
non-Archimedean slope.  This applies simultaneously to energy, \(J\), and
the affine translation term.  It does not alter the DH measure, whose
coordinate \(\lambda\) was fixed by the filtration before a complex norm
convention was chosen.

\paragraph{Firewall boundary.}
Statements quoted verbatim from \cite{BHJ19} retain that paper's ordinary
norm-weight convention, as in
\Cref{thm:rf-bhj-unaffected-slopes}.  Statements formulated in the book's
squared-norm convention use \(t=2s\) and the quotients in
\eqref{eq:rf-factor-two-slope-quotients}.  No line below converts one
convention by changing only the numerator or only the radial parameter.

\section{Algebraic test configurations and the Phong--Sturm ray}

Let $X$ be a connected projective complex manifold of dimension $n$ and let the ample line bundle $L$ carry the smooth positive metric $h$. Use the fixed smooth probability volume form $\mathrm dV$ and the $L^2$-norms of \Cref{sec:rf-filtered-test-curves}. For a decreasing, left-continuous, linearly bounded multiplicative filtration $\mathcal F$, choose an $L^2$-orthonormal basis $s_{m,1},\ldots,s_{m,N_m}$ adapted to the filtration, with weights $\lambda_{m,j}=\chi_m(s_{m,j})$. The weighted Bergman potential is
\[
 \Phi_{m,t}:=\frac1m\log\left(
       \sum_{j=1}^{N_m}\rme^{t\lambda_{m,j}}
                    |s_{m,j}|_{h^m}^2\right),\qquad t\geq0.
\tag{65.23}\label{eq:rf-bergman-ray}
\]
It is independent of the choice of orthonormal bases in the orthogonal graded pieces. Degrees are sufficiently divisible for a rational polarization. The exponential uses the squared-norm radial variable $t=-\log|\zeta|^2$.

\begin{lemma}\label{lem:rf-weighted-bergman-comparison}
For all large admissible $m$, put
\[
 H_{m,t}:=\sup_\tau(B_{m,\tau}+t\tau),
 \qquad H^{(2)}_{m,t}:=\sup_\tau(B^{(2)}_{m,\tau}+t\tau).
\]
With $\varepsilon_m$ from \eqref{eq:rf-bergman-norm-comparison},
\[
 \begin{gathered}
 H^{(2)}_{m,t}\leq\Phi_{m,t}
           \leq H^{(2)}_{m,t}+\frac{\log N_m}{m},\\
 H_{m,t}\leq\Phi_{m,t}
           \leq H_{m,t}+\varepsilon_m+\frac{\log N_m}{m}.
 \end{gathered}
\tag{65.23a}\label{eq:rf-weighted-bergman-comparison}
\]
The bounds are uniform in $x\in X$, $t\geq0$ and the filtration.
\end{lemma}

\begin{proof}
Write $a_j=|s_{m,j}(x)|_{h^m}^2$. For every threshold $\tau$,
\[
 \sum_j\rme^{t\lambda_{m,j}}a_j
   \geq\rme^{mt\tau}\sum_{\lambda_{m,j}\geq m\tau}a_j,
\]
since $t\geq0$. This proves the lower bound by $H^{(2)}_{m,t}$. Conversely the summand with index $j$ is at most $\exp(mH^{(2)}_{m,t})$, by taking threshold $\lambda_{m,j}/m$ and using the partial sum containing $a_j$. Summing the $N_m$ bounds proves the upper bound. Global generation makes all three full suprema finite in large degrees. The second line follows from \Cref{lem:rf-bergman-norm-comparison}. The logarithmic sum estimate compares $\Phi_{m,t}$ directly with the $L^2$ obstacles; passage to the sup-norm obstacles costs the additional $\varepsilon_m$.
\end{proof}

The norm error cannot in general be omitted, even in large degrees. On $(\mathbb P^1,\mathcal O(1))$ with the Fubini--Study metric, put $p=|x|^2/(|x|^2+|y|^2)$ and choose $\mathrm dV=(p+1/2)\omega_{\mathrm{FS}}$, with $\int\omega_{\mathrm{FS}}=1$. This is a smooth positive probability volume form. Give all sections weight zero. Monomials are orthogonal for this volume form, and
\[
 \|y^m\|_{2,m}^2
   =\int_0^1(1-p)^m(p+1/2)\,\mathrm dp
   =\frac{m+4}{2(m+1)(m+2)}.
\]
At $p=0$ only $y^m$ contributes to the Bergman sum. Thus $H_{m,t}=0$, while
\[
 \Phi_{m,t}=\frac1m\log\frac{2(m+1)(m+2)}{m+4}
              >\frac{\log(m+1)}m=\frac{\log N_m}m
 \qquad(m\geq1).
\]
The logarithmic sum error alone therefore does not compare the two norm conventions.

\begin{theorem}\label{thm:rf-phong-sturm}
For the filtration above, the usc tail envelope
\[
 \Phi_t:=\lim_{k\to\infty}
               \left(\sup_{m\geq k}\Phi_{m,t}\right)^*
\tag{65.24}\label{eq:rf-bergman-limit}
\]
equals $\ell_t^{\mathcal F}$ for $t\geq0$. For a normal relatively ample test configuration, with filtration \eqref{eq:rf-tc-filtration}, it is the Phong--Sturm weak geodesic ray in the fixed-linearization convention of \cite[Definition~9.1]{RWN14}. The limiting ray is independent of the auxiliary smooth probability volume form.
\end{theorem}

\begin{proof}
Put $\delta_m=\varepsilon_m+m^{-1}\log N_m$. Since $N_m=O(m^n)$, one has $\delta_m\to0$. The pointwise inequality $B_{m,\tau}\leq g_\tau^{\mathcal F}$ and the equality of the raw and closed Legendre suprema give
\[
 \Phi_{m,t}\leq\ell_t^{\mathcal F}+\delta_m.
\]
The ray is already psh and upper semicontinuous, so upper regularization preserves this bound. Taking a tail and then its decreasing limit gives $\Phi_t\leq\ell_t^{\mathcal F}$.

For the reverse inequality, for every fixed $\tau$ the partial-sum estimate in the preceding proof gives
\[
            \Phi_{m,t}\geq B^{(2)}_{m,\tau}+t\tau.
\]
Take tail suprema and upper regularizations, then use \eqref{eq:rf-partial-bergman-tail} to obtain $\Phi_t\geq g_\tau^{\mathcal F}+t\tau$. Taking the supremum over $\tau<c_{\mathcal F}$ recovers $\ell_t^{\mathcal F}$, including the closed endpoint by its left limit. Thus both inequalities are proved without interchanging a parameter supremum with a limit. At $t=0$ the same bounds give $\Phi_0=0$. The estimates also apply to the joint envelopes on compact time intervals.

For a test configuration, the resulting construction is \cite[Definition~9.1 and Theorem~9.2]{RWN14}, with the original weak-ray existence theorem in \cite[Theorem~1.1]{PS07}. The two inequalities above reconstruct the published proof of Theorem~9.2 using the uniform estimate from its Lemma~7.6. Since the sup-norm obstacles and their test curve do not involve $\mathrm dV$, equality with their ray proves the last assertion.
\end{proof}

The fixed-linearization convention retains all weights. The original construction uses the traceless part of the generator and also rescales the basis by $m^{-n/2}$; see \cite[Sections~2.1 and~4.2]{PS07}. In the book's weight sign and time convention, put
\[
 \overline\lambda_m=\frac{1}{mN_m}\sum_j\lambda_{m,j}.
\]
Replacing $\lambda_{m,j}$ by $\lambda_{m,j}-m\overline\lambda_m$ changes $\Phi_{m,t}$ exactly to $\Phi_{m,t}-t\overline\lambda_m$. By the compactly supported spectral convergence in \Cref{thm:rf-dh-quantization}, $\overline\lambda_m$ tends to $\overline\lambda=\int\lambda\,\mathrm d\DHm^{\mathrm{prob}}(\mathcal F)$. The centered tail envelope is consequently $\ell_t^{\mathcal F}-t\overline\lambda$; the error $t|\overline\lambda_m-\overline\lambda|$ is uniform on compact time intervals. The basis rescaling subtracts $n\log m/m$ and has zero limit. These scalar corrections must be distinguished from the radial factor of two in \Cref{sec:rf-factor-two-firewall}.

The conclusion is equality of usc tail envelopes. No smooth convergence or convergence of curvature tensors is inferred from the logarithmic sum estimate.

\section{The exact commutative diagrams}

There are two useful diagrams for finitely generated rational
filtrations.
The first is algebraic and uses the relatively ample Rees configuration before normalization, together with its section-saturated filtration; the Veronese index is chosen so that this is the given filtration. It lives over the trivially valued field:
\[
\begin{array}{ccc}
\left.
 \begin{array}{c}
 \text{relatively ample}\\[-1mm]
 (\mathcal X,\mathcal L)
 \end{array}
\right\}
&\xrightarrow{\ \eqref{eq:rf-tc-filtration}\ }&
\mathcal F
\\[3mm]
\downarrow&&\downarrow
\\[-1mm]
\phi_{(\mathcal X,\mathcal L)}
&=&
\phi_{\mathcal F}.
\end{array}
\tag{65.25}\label{eq:rf-algebraic-square}
\]
The lower equality is an equality of rational PL metrics on
\(X^{\An}\), including at proper-support semivaluations.

The second diagram requires the projective complex hypotheses of
\Cref{chap:dxz-bj-comparison}:
\[
\begin{array}{ccccc}
\mathcal F
&\longrightarrow&
\Gamma^{\mathcal F}
&\xleftrightarrow{\ \mathrm{Legendre}\ }&
\ell^{\mathcal F}
\\[2mm]
\downarrow&&\downarrow\mathfrak A&&\Vert
\\[-1mm]
\phi_{\mathcal F}
&=&
(\Gamma^{\mathcal F})^{\An}
&&\ell^{\mathcal F}.
\end{array}
\tag{65.26}\label{eq:rf-analytic-square}
\]

\begin{theorem}\label{thm:rf-commutativity}
Diagram \eqref{eq:rf-algebraic-square} commutes for every finitely generated
rational filtration.
Diagram \eqref{eq:rf-analytic-square} commutes for such a filtration on a
projective complex manifold with an ample \(\mathbf Q\)-line bundle.
For a test-configuration filtration, its right-hand ray is the
Phong--Sturm ray.
\end{theorem}

\begin{proof}
For \eqref{eq:rf-algebraic-square}, take homogeneous Rees generators.
Both lower paths give the finite maximum \eqref{eq:rf-finite-fs}. Passing to the normalization changes neither lower function by \Cref{lem:rf-integral-closure}; the upper filtration can change, as \Cref{ex:rf-normalization-changes-weights} shows.

For \eqref{eq:rf-analytic-square}, the valuative analytification of the
partial Bergman level records the base ideal of
\(\mathcal F^{m\tau}R_m\).
Taking the supremum over \(m\), then the maximal envelope, gives
\[
 (\Gamma^{\mathcal F})^{\An}(v)
 =\sup_{m,s\ne0}\frac{\chi_m(s)-v(s)}m
 =\phi_{\mathcal F}(v)
\]
first on divisorial valuations.
Both sides are Boucksom--Jonsson psh metrics, so divisorial determination
\cite[Theorem~4.22]{BJ22} extends the equality from \(X^{\divis}\) to all of
\(X^{\An}\).  That already covers the proper-support strata, at which
neither side need be \(-\infty\): by \eqref{eq:rf-linear-bounds} the base
ideal \(\mathfrak a_{m,mC_-}\) is trivial for large \(m\), so both sides are
bounded below by \(C_-\) everywhere, in agreement with
\Cref{prop:rf-fg-pl}, which asserts that both are rational PL and hence real
valued.  (\Cref{ex:rf-p1-complete} makes this concrete: at the trivial
valuation supported at \(p=[0:1]\) both sides equal \(0\), not
\(-\infty\), since \(y^m\) does not vanish at \(p\).)  No separate argument
is needed there.  This step is not formal --- upper semicontinuity alone gives
inequalities in only one direction --- and it is exactly what licenses the
claim that the lower equality holds at proper-support semivaluations.
The right square is the Ross--Witt Nystr\"om correspondence.
The final assertion is \Cref{thm:rf-phong-sturm}.
\end{proof}

\begin{remark}[Scope of the vertical equality]
The equality involving \(\mathfrak A\) uses the projective DXZ--BJ
comparison.
It is not asserted for an arbitrary compact K\"ahler class.
The proof just given does not pass through Mesquita--Piccione theory.  In the
connected smooth projective ample \(\mathbf Q\)-class setting, however,
\Cref{thm:pseff-projective-full-mp-bj} separately identifies the full BJ and
MP psh cones; this does not extend the displayed equality to the general
compact K\"ahler or non-nef settings.
\end{remark}

\section{Duistermaat--Heckman quantization}

Let
\(\lambda_{m,1}\geq\cdots\geq\lambda_{m,N_m}\)
be the filtration jumping numbers and set
\[
 \nu_m=\frac1{N_m}\sum_{j=1}^{N_m}
                 \delta_{\lambda_{m,j}/m}.
\tag{65.27}\label{eq:rf-spectral-measure}
\]

\begin{theorem}\label{thm:rf-dh-quantization}
For a linearly bounded multiplicative filtration of the ample section ring,
the probability measures \(\nu_m\) converge weakly, along sufficiently
divisible \(m\), to a compactly supported probability measure
\(\DHm^{\mathrm{prob}}(\mathcal F)\).
For a test-configuration filtration this is the BHJ
Duistermaat--Heckman measure.
The unnormalized measure of the associated test curve is
\[
                    \DHm(\Gamma^{\mathcal F})
                    =V\DHm^{\mathrm{prob}}(\mathcal F).
\tag{65.28}\label{eq:rf-dh-mass-conversion}
\]
\end{theorem}

\begin{proof}
Linear boundedness puts all supports in a fixed compact interval.
For a continuity point \(\tau\), the tail distribution of \(\nu_m\) is
\[
 \nu_m([\tau,+\infty))
 =\frac{\dim\mathcal F^{m\tau}R_m}{N_m}.
\]
The volume theorem for the multiplicative filtered subseries gives the
limit of the right side; this is the filtered Okounkov-body limit
\cite[Theorem~A]{BC11}.
Compactness upgrades convergence of distribution functions to weak
convergence.
For test configurations this is
\cite[Corollary~3.4]{BHJ17}.

The filtered Okounkov body has Euclidean volume \(V/n!\).
Its concave transform has the same tail distribution, and the book's factor
\(n!\) makes its pushforward have mass \(V\).
This proves \eqref{eq:rf-dh-mass-conversion}.
\end{proof}

The Bergman rays \eqref{eq:rf-bergman-ray} and spectral measures
\eqref{eq:rf-spectral-measure} are two quantizations of the same filtration.
The first converges to a weak geodesic ray, the second to its asymptotic
slope distribution.
Neither convergence alone implies convergence of curvature tensors.

\subsection{Finite-level moment checks}

For the action in \Cref{ex:rf-p1-complete}, the degree-\(m\) weights are
\(0,1,\ldots,m\).  Before taking a limit one already has
\[
 \int\lambda\,\rd\nu_m
       =\frac1{m+1}\sum_{j=0}^m\frac jm=\frac12,
 \qquad
 \int\lambda^2\,\rd\nu_m
       =\frac{2m+1}{6m}\longrightarrow\frac13.
\tag{65.29}\label{eq:rf-p1-finite-moments}
\]
The first equality is exact at every level; the second records the first
nontrivial convergence error.  Since the supports remain in \([0,1]\),
these moment limits agree with Lebesgue measure on that interval.

For the complex ray, use the Fubini--Study metric on $\mathcal O(1)$ and its curvature volume form of mass one. The squared $L^2$-norm of $x^j y^{m-j}$ is
\[
 \int_0^1p^j(1-p)^{m-j}\,\mathrm dp
      =\frac{j!(m-j)!}{(m+1)!},
\]
where $p=|x|^2/(|x|^2+|y|^2)$. Angular integration makes distinct monomials orthogonal, so the adapted orthonormal basis is
\[
                 s_{m,j}=\sqrt{(m+1)\binom mj}\,x^jy^{m-j}.
\]
The binomial theorem gives the exact identity
\[
 \Phi_{m,t}
   =\frac{\log(m+1)}m+
       \log\frac{|y|^2+\rme^t|x|^2}{|x|^2+|y|^2}.
\tag{65.30}\label{eq:rf-p1-log-sum-limit}
\]
Thus the ray is $\ell_t=\log(1-p+\rme^t p)$, and the error is the constant $m^{-1}\log(m+1)$ for every $x$ and $t\geq0$. Since $\overline\lambda_m=1/2$ in every degree, the centered ray is $\ell_t-t/2$. The orthonormalizing coefficients are essential. On the chart $y\ne0$, the unnormalized sum $\sum_j\rme^{tj}|x/y|^{2j}$ instead gives the limit $\max\{0,t+\log|x/y|^2\}$ before subtracting the reference weight. At $t=0$ this would give $\max\{0,\log|x/y|^2\}-\log(1+|x/y|^2)$, which is not the chosen zero reference. It is therefore not the Bergman quantization in \eqref{eq:rf-bergman-ray}.

For the induced action on \(\mathcal O(d)\), the degree-\(m\) weights are
\(0,1,\ldots,dm\).  The probability DH limit has density \(d^{-1}\) on
\([0,d]\), whereas the unnormalized test-curve measure has mass
\(V=d\) and is ordinary Lebesgue measure on \([0,d]\).  Thus
\[
 \int\lambda\,\rd\DHm^{\mathrm{prob}}=\frac d2,
 \qquad
 \int\lambda\,\rd\DHm=\frac{d^2}{2}.
\tag{65.31}\label{eq:rf-od-mass-moment-check}
\]
The example separates the support scale \(d\) from the total-mass factor
\(V=d\); they happen to be numerically equal here but enter for different
reasons.

\section{Radial energy slopes}

Normalize the complex Monge--Amp\`ere energy by its mass, so
\(E(\varphi+c)=E(\varphi)+c\).
Let \(E^{\mathrm u}=VE\) denote the unnormalized energy.

\begin{theorem}\label{thm:rf-energy-slope}
Let \(\mathcal F\) be a linearly bounded filtration whose test curve has
finite energy, and let \(\ell^{\mathcal F}\) be its Legendre ray.
Then
\[
 \lim_{t\to\infty}\frac{E^{\mathrm u}(\ell_t^{\mathcal F})}{t}
 =\int_{\mathbf R}\lambda\,
          \rd\DHm(\Gamma^{\mathcal F})(\lambda),
\tag{65.32}\label{eq:rf-energy-first-moment}
\]
or equivalently
\[
 \lim_{t\to\infty}\frac{E(\ell_t^{\mathcal F})}{t}
 =\int_{\mathbf R}\lambda\,
          \rd\DHm^{\mathrm{prob}}(\mathcal F)(\lambda).
\tag{65.33}\label{eq:rf-normalized-energy-slope}
\]
For a normal ample test configuration,
\[
 \lim_{t\to\infty}\frac{E^{\mathrm u}(\ell_t)}t
 =\frac{(\overline{\mathcal L}^{n+1})}{n+1}.
\tag{65.34}\label{eq:rf-energy-intersection-slope}
\]
\end{theorem}

\begin{proof}
The Ross--Witt Nystr\"om energy identity identifies the radial energy slope
with the layer-cake energy of the test curve.
The first-moment formula
\eqref{eq:cma-first-moment} identifies that layer-cake energy with the
unnormalized DH first moment.
Divide by \(V\) to obtain
\eqref{eq:rf-normalized-energy-slope}.
For a test configuration, the BHJ intersection-moment formula
\eqref{eq:cma-dh-first-intersection} gives
\eqref{eq:rf-energy-intersection-slope}.
\end{proof}

These energy identities use the RWN correspondence and the unaffected DH
intersection theorem of \cite{BHJ17}.
They do not depend on \cite[Theorem~5.6]{BHJ19}.

\begin{theorem}\label{thm:rf-bhj-unaffected-slopes}
Let \((\mathcal X,\mathcal L)\) be an ample test configuration for a
polarized complex manifold \((X,L)\).
Let \(\Phi\) be a smooth strictly positive \(S^1\)-invariant metric on
\(\mathcal L\) near the central fibre, and let \((\psi_s)\) be the
corresponding smooth positive ray in the ordinary norm-weight convention of
\cite{BHJ19}.
Then the asymptotic formulas of
\cite[Theorem~3.6]{BHJ19} identify the slopes of
\(E,H,I,J,M,R\) with
\(E^{\NA},H^{\NA},I^{\NA},J^{\NA},M^{\NA},R^{\NA}\), respectively.
In the logarithmic setting, the same identifications hold under the
subklt and positivity hypotheses of \cite[Theorem~4.2]{BHJ19}.
\end{theorem}

\begin{proof}
We reconstruct the proof of \cite[Theorem~3.6]{BHJ19}, keeping the
source's norm-weight convention throughout.  Two results proved in that
paper are foundational imported inputs for this reconstruction:

\begin{enumerate}
\item the Deligne-pairing asymptotic formula
      \cite[Lemma~3.9]{BHJ19};
\item the local simple-normal-crossings integral estimates
      \cite[Lemmas~3.10--3.11]{BHJ19}.
\end{enumerate}

Neither input belongs to Section~5.  We now explain how they give every
functional in the statement.

\smallskip
\noindent\emph{Step 1: reduction to a smooth dominating model.}
After an equivariant resolution and equivariant compactification over
\(\mathbf P^1\), we may suppose that the total space is smooth, dominates
\(X\times\mathbf P^1\), and has simple-normal-crossings central fibre.
All six non-Archimedean functionals in the statement are invariant under
this pullback.  On the complex side the induced ray does not change,
because it is read on the common general fibre.  It is therefore enough to
work on this model.

\smallskip
\noindent\emph{Step 2: the Deligne-pairing slope.}
Let \(\mathcal M_0,\ldots,\mathcal M_n\) be equivariant line bundles on
the compactified model, supplied with smooth invariant metrics near the
central fibre, and let \(m_{j,s}\) be their induced metrics on the general
fibre.  Lemma~3.9 of the source gives, in its normalization,
\[
 \lim_{s\to+\infty}\frac1s
 \bigl\langle m_{0,s},\ldots,m_{n,s}\bigr\rangle
 =
 \bigl(\overline{\mathcal M}_0\cdots
       \overline{\mathcal M}_n\bigr).
\tag{65.35}\label{eq:rf-deligne-basic-slope}
\]
Here the left-hand side denotes the logarithmic change of the metric on
the Deligne pairing relative to fixed reference metrics, and the right-hand
side is the compactified intersection number.  To see why the dependence
is linear in \(s\), choose a meromorphic trivialization of the Deligne
pairing over the disk.  Equivariance writes it as a power of the base
parameter times a nowhere-vanishing local frame.  With
\(s=-\log|t|\), the power of \(t\) contributes the linear term, while the
logarithm of the norm of the nonvanishing frame is bounded.  The projection
formula identifies the coefficient with the intersection number in
\eqref{eq:rf-deligne-basic-slope}.  This argument also fixes the sign: the
same equivariant trivialization is used in the definitions of the
non-Archimedean intersections.

The Deligne-pairing expressions for \(E,I,J\), and \(R\), recalled in
\cite[Lemma~1.2]{BHJ19}, are finite linear combinations of terms of the
form appearing in \eqref{eq:rf-deligne-basic-slope}.  Applying
\eqref{eq:rf-deligne-basic-slope} term by term gives
\[
 \lim_{s\to+\infty}\frac{F_{\mathrm{wt}}(\psi_s)}s
     =F^{\NA}(\phi^{\NA}),
 \qquad F\in\{E,I,J,R\}.
\tag{65.36}\label{eq:rf-deligne-functional-slopes}
\]
No interchange of an infinite limit with a Monge--Amp\`ere expansion occurs
here: all pairings are multilinear expressions with exactly \(n+1\)
entries.

\smallskip
\noindent\emph{Step 3: entropy.}
The source first treats \(H+R\), because that sum has a direct
Deligne-pairing expression.  Choose a smooth invariant metric
\(\Psi_{\mathrm{ref}}\) on the logarithmic relative canonical bundle of
the resolved total space, and denote its fibrewise metric by
\(\psi_{\mathrm{ref}}^s\).  The Deligne identity used in the source is
\[
\begin{aligned}
 V\bigl(H(\psi_s)+R(\psi_s)\bigr)
  ={}&
 \left\langle\frac12\log\MA(\psi_s),
       \psi_s,\ldots,\psi_s\right\rangle_X \\
 &-\left\langle\psi_{\mathrm{ref}},
       \psi_{\mathrm{ref},L},\ldots,
       \psi_{\mathrm{ref},L}\right\rangle_X .
\end{aligned}
\tag{65.37}\label{eq:rf-entropy-deligne-decomposition}
\]
The names \(\psi_{\mathrm{ref}}\) and
\(\psi_{\mathrm{ref},L}\) in the second line denote fixed reference
metrics on \(K_X\) and \(L\), respectively.  Replacing the first entry in
the first pairing by \(\psi_{\mathrm{ref}}^s\),
\eqref{eq:rf-deligne-basic-slope} gives the intersection
\[
       \bigl(K^{\log}_{\overline{\mathcal X}/\mathbf P^1}
              \cdot\overline{\mathcal L}^{n}\bigr).
\]
The error made by this replacement is exactly
\[
 \frac12\int_X
   \log\!\left[
      \frac{\MA(\psi_s)}
           {\rme^{2\psi_{\mathrm{ref}}^s}}
   \right](\ddc\psi_s)^n .
\tag{65.38}\label{eq:rf-entropy-replacement-error}
\]

We spell out why \eqref{eq:rf-entropy-replacement-error} is \(o(s)\).
On an snc chart with
\(t=u z_0^{b_0}\cdots z_p^{b_p}\), logarithmic polar coordinates identify
the relevant part of the fibre with a fixed simplex, a compact angular
group, and a transverse polydisc.  The logarithmic canonical volume form
then has mass comparable to \(s^p\).  After summing finitely many charts,
\cite[Lemma~3.11]{BHJ19} gives
\[
 C^{-1}s^d\leq
       \int_{X_t}\rme^{2\psi_{\mathrm{ref}}^s}
       \leq Cs^d,
 \qquad 0\leq d\leq n,
\tag{65.39}\label{eq:rf-entropy-polynomial-bound}
\]
where \(d\) is the dimension of the dual complex of the snc central
fibre.  Normalize this volume form to a probability measure.  Positivity
of relative entropy and \eqref{eq:rf-entropy-polynomial-bound} give the
lower bound \(-O(\log s)\) for twice
\eqref{eq:rf-entropy-replacement-error}.

For the upper bound, put
\[
 g_t:=
 \frac{(\ddc\Phi|_{X_t})^n}
      {\rme^{2\Psi_{\mathrm{ref}}|_{X_t}}}.
\]
Smooth semipositivity gives \(\ddc\Phi\leq C\Omega\) on each of the
finitely many charts.  The same logarithmic-coordinate calculation shows
that \((\Omega|_{X_t})^n/
\rme^{2\Psi_{\mathrm{ref}}|_{X_t}}\) is uniformly bounded.  Hence \(g_t\)
is uniformly bounded above.  Since \((\ddc\Phi|_{X_t})^n\) has fixed mass
\(V\), the integral in \eqref{eq:rf-entropy-replacement-error} is uniformly
bounded above.  Together with the lower bound, this is exactly the
\(O(\log s)\) estimate of \cite[Lemma~3.10]{BHJ19}.

It follows from \eqref{eq:rf-entropy-deligne-decomposition} that the slope
of \(H+R\) equals the non-Archimedean slope of \(H^{\NA}+R^{\NA}\).
Subtracting the already established \(R\)-identity in
\eqref{eq:rf-deligne-functional-slopes} yields
\[
       \lim_{s\to+\infty}\frac{H_{\mathrm{wt}}(\psi_s)}s
          =H^{\NA}(\phi^{\NA}).
\tag{65.40}\label{eq:rf-entropy-smooth-slope}
\]
In the logarithmic version, the subklt hypothesis is what makes the local
power singularities integrable.  The analogous conclusion without that
hypothesis is not asserted here.

\smallskip
\noindent\emph{Step 4: Mabuchi energy.}
The Chen--Tian decomposition in the source convention expresses \(M\) as
the sum of \(H\), \(R\), and the appropriate scalar multiple of \(E\).
Combining \eqref{eq:rf-deligne-functional-slopes} and
\eqref{eq:rf-entropy-smooth-slope}, with the identical decomposition on
the non-Archimedean side, proves the formula for \(M\).  This completes the
six slope calculations.

Finally, the authors' erratum \cite{BHJ19err} explicitly states that
Theorem~A, Corollary~B, and all results through Section~4 are independent
of the flawed Theorem~5.6.  The construction above uses only material from
Sections~1--4.  It therefore gives the claimed formulas without invoking
any quarantined statement.
\end{proof}

\Cref{thm:rf-bhj-unaffected-slopes} concerns the smooth ray induced
by a smooth metric near the central fibre.
Transporting every functional slope to an arbitrary weak finite-energy ray
requires an additional extension theorem for that functional.
Only the energy extension proved in
\Cref{thm:rf-energy-slope} is used here without further hypotheses.

\section{The BHJ 2019 erratum quarantine}

The dependency graph used in this chapter is
\[
\begin{array}{ccc}
\cite{BHJ17}\text{ DH/intersections}
 &\longrightarrow&
\text{energy moments and slopes}
\\[1mm]
\cite{RWN14}\text{ Legendre energy}
 &\longrightarrow&
\text{energy moments and slopes}
\\[1mm]
\begin{array}{c}
\cite[Sections~3--4]{BHJ19}\\[-1mm]
\text{unaffected}
\end{array}
 &\longrightarrow&
\text{smooth-ray asymptotics}
\\[1mm]
\cite[Theorem~5.6]{BHJ19}
 &\dashrightarrow&
\begin{array}{c}
\text{Theorem~5.4 and Theorem~C}\\[-1mm]
\text{Corollaries~D--E}
\end{array}.
\end{array}
\tag{65.41}\label{eq:rf-erratum-dependency-graph}
\]
The dashed bottom row is quarantined.  No path in either solid row passes
through it.

The precise unsafe zone is short enough to record verbatim as a dependency
ledger.

\begin{theorem}\label{thm:rf-erratum-audit}
The following results of \cite{BHJ19} are not used as proved statements in
this book:
\[
 \text{Theorem~5.6},\qquad
 \text{Theorem~5.4},\qquad
 \text{Theorem~C},\qquad
 \text{Corollaries~D and~E}.
\tag{65.42}\label{eq:rf-quarantined-results}
\]
The proof of Theorem~5.6 contains an incorrect claim.
The published erratum says that its statement may still be true, but is
unproved by that argument.
Theorem~5.4 and the last three results depend on it.

Theorem~A, Corollary~B, and every result through Section~4 remain valid
according to the erratum.
\end{theorem}

\begin{proof}
This is exactly the dependency statement in Section~2 of
\cite{BHJ19err}.
The erratum identifies the false proof step: for a one-parameter subgroup
\(\lambda\) and a \(K\)-point \(\xi\), the reduction of
\(\lambda\cdot\xi\) need not depend only on the reduction of \(\xi\).
Consequently the proposed orbit-closure argument does not establish
Theorem~5.6.
The same erratum lists Theorem~5.4, Theorem~C, and Corollaries~D--E as
affected, and separately confirms the unaffected results stated above.
\end{proof}

\begin{warning}{No surrogate use}
\label{warn:rf-no-surrogate}
We do not use a reformulation of a quarantined result under another name.
In particular, no implication from uniform K-stability to coercivity on all
Bergman levels is deduced from Theorem~C or Corollaries~D--E of
\cite{BHJ19}.
Any later theorem with a similar conclusion must carry an independent
post-erratum source and its own hypotheses.
\end{warning}

\section{A complete example on \texorpdfstring{\(\mathbf P^1\)}{P1}}

\begin{example}
\label{ex:rf-p1-complete}
Let \(X=\mathbf P^1\), \(L=\mathcal O(1)\), and let
\(a\cdot[x:y]=[ax:y]\).
Choose the linearization for which
\(x^jy^{m-j}\) has weight \(j\).
Then
\[
 \mathcal F^\lambda R_m
 =\operatorname{span}\{x^jy^{m-j}\mid j\geq\lambda\}.
\tag{65.43}\label{eq:rf-p1-filtration}
\]
The normalized spectral measures are
\[
 \nu_m=\frac1{m+1}\sum_{j=0}^m\delta_{j/m}
        \rightharpoonup\mathbf 1_{[0,1]}\,\rd \lambda.
\tag{65.44}\label{eq:rf-p1-dh}
\]
Since \(V=1\), probability and unnormalized DH measures coincide.
The normalized energy slope is
\[
                 \int_0^1\lambda\,\rd \lambda=\frac12.
\tag{65.45}\label{eq:rf-p1-energy-slope}
\]

Use the Fubini--Study reference and write $p=|x|^2/(|x|^2+|y|^2)$. Formula \eqref{eq:rf-p1-log-sum-limit} gives the ray $\ell_t=\log(1-p+\rme^t p)$. For $0<p<1$, minimizing $\ell_t-t\tau$ over $t>0$ gives
\[
 \Gamma_\tau^{\mathcal F}
 =\begin{cases}
  0,&\tau\leq p,\\
  \displaystyle\tau\log\frac p\tau
       +(1-\tau)\log\frac{1-p}{1-\tau},&p<\tau<1,\\
  \log p,&\tau=1,\\
  -\infty,&\tau>1.
 \end{cases}
\]
Indeed, for $p<\tau<1$ the minimizing time satisfies $\rme^t=\tau(1-p)/(p(1-\tau))$; for $\tau\leq p$ the infimum is the limit at $t=0$. At $p=0$ the level is zero for $\tau\leq0$ and $-\infty$ for $\tau>0$; at $p=1$ it is zero for $\tau\leq1$ and $-\infty$ for $\tau>1$. These are also the psh limiting values of the displayed formula. The endpoint is $\log p$ and the initial levels are zero for $\tau\leq0$, in agreement with the hyperplane calculation in \Cref{ex:rays-hyperplane-test-curve}.

Twisting the linearization by \(c\) changes every weight to \(j+mc\).
Equations \eqref{eq:rf-translation} and
\eqref{eq:rf-p1-dh} then give
\[
 \ell_t\mapsto\ell_t+ct,
 \qquad
 \DHm\mapsto(\lambda\mapsto\lambda+c)_*\DHm,
 \qquad
 \frac12\mapsto\frac12+c.
\]
This simultaneously checks the Legendre sign, the DH translation, and the
energy normalization.
\end{example}

\section{What the diagram does not say}

The construction diagram has four strict boundaries.

\begin{enumerate}
\item A general linearly bounded filtration defines a ray, but need not
define a finite-type test configuration.

\item In a merely big class, \cite[Proposition~7.11]{RWN14} alone does not
remove the outer maximal envelope in
\eqref{eq:rf-filtration-test-curve}.  It is removed instead by
\Cref{thm:pseff-filtration-maximality}, which needs no finite-generation
hypothesis, and only on the strict interior \(\tau<\Gamma_{\max}\).

\item Weak convergence of Bergman rays and spectral measures does not imply
smooth convergence or convergence of all energy functionals.

\item The projective DXZ--BJ vertical arrow in
\eqref{eq:rf-analytic-square} is not an MP comparison theorem and is not
extended to an arbitrary compact K\"ahler class here.
\end{enumerate}

\section*{Source note}

The filtration, Rees algebra, and trivially valued FS dictionary are
\cite[Sections~1--2]{BJ22} and \Cref{chap:bj-test-configurations}.
The algebraic DH limit and intersection moments are \cite{BHJ17}.
Analytic test curves, their maximal envelopes, filtration rays, and the
identification with the Phong--Sturm ray are
\cite[Theorems~1.1, 1.2, and~9.2]{RWN14}; the original Bergman-ray
construction is \cite{PS07}.
The projective analytification arrow is \cite{DXZ23} in the exact domain of
\Cref{chap:dxz-bj-comparison}.
The asymptotic functional theorem used above is the unaffected
\cite[Theorem~3.6]{BHJ19}.
The quarantine \eqref{eq:rf-quarantined-results} follows
\cite{BHJ19err} and is a hard boundary for every later slope argument.
All Deligne-pairing identities used in this chapter come from the cited
unaffected Sections~1--4 of \cite{BHJ19} or from the already established
intersection formulas.  No excluded or derivative source is used.
